\documentclass[a4paper, 11pt]{article}

% -- Coments
\usepackage{verbatim}

% ----- Fonts -----
% -- Fuente --
% \usepackage{fontspec}
% \setmonofont{JetBrainsMono Nerd Font}

\usepackage{babel}
\babelprovide[import, main]{spanish}

% -- Color --
\usepackage{xcolor}
%\definecolor{azul}{RGB}{00,33,99}
\definecolor{azul}{RGB}{35,72,180}

% -- Page Margin --
\usepackage[margin=1in, headheight=15pt]{geometry}

% -- Encabezado con nombre y expediente --
\usepackage{fancyhdr}
\pagestyle{fancy}
\fancyhf{}
\lhead{\textbf{Bórquez Guerrero Angel Fernando}}
\rhead{Expediente: 219208106}
\renewcommand{\headrulewidth}{0.5pt}

% -- Espaciados --
\newcommand{\Pspace}{0.5cm}
\newcommand{\Aspace}{0.2cm}

% -- Columnas --
\usepackage{multicol}

% -- Imagenes --
\usepackage{graphicx}
\usepackage{float}

% -- Matemáticas --
\usepackage{amsmath, amssymb}

% -- Gráficas --
\usepackage{pgfplots}
\pgfplotsset{compat=1.18}


\title
{
    \textbf{Contexto Regional 2026-2} \\[0.3cm]
    {\large \textit{Resistencias por debajo del agua}} \\[0.2cm]
    {\normalsize Celaya Aguilar, Suzette} \\[0.3cm]
    {\normalsize Profesor José Moreno Vega}
}
\date{}

\begin{document}

\begin{otherlanguage}{spanish}
\maketitle
\end{otherlanguage}
\thispagestyle{fancy}
\vspace{-0.5cm}


{\Large 1. ¿Cuál es el argumento central del autor?} \\
{\color{azul}
    La autora sostiene que los desplazados por la presa El Novillo, que en 1964 inundó los pueblos de Batuc, Tepupa y Suaqui, no fueron simples víctimas ni un dato estadístico, sino que resistieron a su manera. El plan de indemnización de 1962 se basaba en la propiedad comprobable y en el ganado, así que dejaba en desventaja a quienes no tenían títulos, y además se presentó con muy poco tiempo para rechazarlo. Frente a eso, algunos desplazados fundaron nuevos poblados y para lograrlo solicitaron tierras ejidales entre 1960 y los primeros años de 1990. Lo importante es que no recurrieron a protestas abiertas contra la CFE o el Estado, sino a un trámite agrario que los tres pueblos ya conocían desde hacía más de un siglo para defender sus tierras. Por eso, apoyándose en Scott, concluye que fundar esos poblados fue una resistencia cotidiana: no desafiaba directamente al poder, pero permitió que el grupo mantuviera su unidad social y productiva y conservara el nombre de los pueblos inundados.
}
\vspace{1cm}

{\Large 2. ¿Qué evidencias utiliza para respaldarlo?} \par
{\color{azul}
Usa sobre todo documentos de archivo, como los del Archivo General Agrario y del Archivo General del Estado de Sonora, donde está el plan de indemnización de 1962 con sus montos por hectárea, por solar y por construcción, además de las opciones de reubicarse en el distrito de riego del Mayo o recibir el pago en efectivo. También presenta el historial de solicitudes de tierras de los tres pueblos desde 1922, como las de Suaqui y Tepupa, y los títulos de la época de 1834 y 1835, para mostrar que ya habían usado el trámite agrario para resolver conflictos. Se apoya en testimonios de los afectados, que sentían que el comité de defensa no los representaba, y en el censo de 1950, el trabajo de Terán y su propio libro \textit{Lo que El Novillo se llevó}. Por último, utiliza teoría de autores como Scott, Bonfil Batalla y Da Rocha.
}
\vspace{1cm}

{\Large 3. ¿Cómo se relaciona esta lectura con la lectura anterior?} \par
{\color{azul}
Las dos hablan de grupos subordinados que consiguen actuar por cuenta propia sin enfrentarse abiertamente al poder. En la lectura 9, las maestras y sindicalistas se movían dentro de los cauces de los gobiernos para ganar espacios, y aquí los desplazados usan la solicitud de tierras ejidales, un mecanismo del propio Estado, para permanecer como grupo. Además, en ambas la relación con el gobierno estatal es clave, y la presa lleva el nombre de Plutarco Elías Calles.
}

\end{document}
